\documentclass[11pt]{article}
\usepackage[margin=1in]{geometry}
\usepackage{amsmath,amssymb,amsthm}
\usepackage{graphicx}
\usepackage{booktabs}
\usepackage{xcolor}
\usepackage{listings}
\usepackage[numbers]{natbib}
\usepackage{hyperref}

\newtheorem{theorem}{Theorem}
\newtheorem{lemma}[theorem]{Lemma}
\theoremstyle{definition}
\newtheorem{definition}[theorem]{Definition}
\theoremstyle{remark}
\newtheorem{remark}[theorem]{Remark}

\lstdefinelanguage{lean}{
  keywords={theorem,lemma,def,abbrev,namespace,end,open,import,by,intro,
    rintro,obtain,have,set,with,fun,at,if,then,else,show,exact,refine,apply,
    rw,simp,simpa,unfold,cases,rcases,by_cases,omega,ring,ext,left,right,constructor},
  sensitive=true,
  comment=[l]{--},
  morecomment=[s]{/-}{-/},
  string=[b]"
}
\lstset{language=lean,basicstyle=\ttfamily\footnotesize,breaklines=true,
  columns=fullflexible,keepspaces=true,
  keywordstyle=\color{blue!60!black}\bfseries,commentstyle=\color{green!40!black},
  frame=single,
  literate=
    {·}{{$\cdot$}}1 {×}{{$\times$}}1 {₁}{{$_1$}}1 {₂}{{$_2$}}1
    {ℕ}{{$\mathbb{N}$}}1 {ℤ}{{$\mathbb{Z}$}}1 {←}{{$\leftarrow$}}1
    {→}{{$\to$}}1 {↔}{{$\leftrightarrow$}}1 {∀}{{$\forall$}}1
    {∃}{{$\exists$}}1 {∈}{{$\in$}}1 {∑}{{$\sum$}}1 {∨}{{$\lor$}}1
    {⊢}{{$\vdash$}}1 {⟨}{{$\langle$}}1 {⟩}{{$\rangle$}}1
    {¬}{{$\neg$}}1 {≠}{{$\neq$}}1 {α}{{$\alpha$}}1 {ᶜ}{{$^{c}$}}1 {∉}{{$\notin$}}1 {∧}{{$\land$}}1 {≤}{{$\le$}}1
}

\title{A machine-checked raster bridge for \emph{Trees Cannot Cut}:\\
enclosed 4-components equal $C_8-\chi$ on every finite raster}
\author{Justin Hart\\Viridis LLC}
\date{10 November 2026}

\begin{document}
\maketitle

\begin{abstract}
For a finite set $B\subset\mathbb{Z}^2$ of barrier pixels, read with 8-connected
barriers and a 4-connected complement, we give a Lean~4 proof of the raster bridge
\[
  \mathrm{enc}(B) \;=\; C_8(B)-\chi(K(B)),
\]
where $\mathrm{enc}(B)$ is the number of finite 4-components of $\mathbb{Z}^2\setminus B$,
$C_8(B)$ is the number of 8-components of $B$, and $\chi(K(B))=V-E+F$ is the Euler
characteristic of the cubical complex formed by the closed pixel squares. The proof
targets the sealed statement \texttt{P1.Bridge} (rung L3 of the project's ladder),
which also asserts that only finitely many components are enclosed. The proof is the
end of a three-rung ladder. L1 is Gray's bit-quad formula for $\chi$. L2
(\texttt{P1.Insertion}) states that inserting one pixel leaves
$\mathrm{enc}-C_8+\chi$ unchanged. We proved it in earlier sessions by an induction in
raster-scan order, a crossing-parity argument and a crack-boundary permutation. L3
follows from L2 by induction from the empty raster, together with a bounding-box
argument for finiteness. In the same file we also give a Lean proof that
$\mathbb{Z}^2\setminus B$ has exactly one infinite 4-component, so the total number of
compartments is $C_8(B)-\chi(K(B))+1$. We explain which assumption of \emph{Trees Cannot Cut}
this discharges and which it does not. The hypothesis \texttt{FaceCount} in that paper's
Lean core says that compartments equal circuit rank plus one for an abstract corridor
graph, and it rests on Alexander duality. On rasters, our result replaces the duality step
by a theorem with no assumptions, with $C_8-\chi$ playing the role of the first Betti number.
It does not discharge \texttt{FaceCount} as stated. Three links are not formalised: a
correspondence between corridor graphs and rasters under which the graph's circuit rank equals $C_8-\chi$
(the solid $2\times2$ block shows that the pixel-adjacency graph is the wrong
choice), the identification of $C_8-\chi$ with $\operatorname{rank}H_1(K(B))$, and the
anchoring at the landscape boundary $\partial\Omega$. The development uses no planar-topology library.
\end{abstract}

\section{Introduction}

Hart's \emph{Trees Cannot Cut}~\cite{Hart2026} argues that fuel breaks and
habitat corridors split a forest landscape into separate compartments only
through their cycles: an acyclic corridor leaves the fuel matrix connected. The
graph theory behind this claim is machine-checked, but the step from topology to
graph theory is carried as an explicit hypothesis (\texttt{FaceCount}), which
says that the number of compartments equals the circuit rank plus one. The
manuscript justifies it by Alexander duality. Mathlib~\cite{Mathlib2020} has
neither Alexander duality nor a planar Euler formula in a form that applies
here.

Fire and fuel-break data usually come as rasters, so project P1 aims to
discharge the hypothesis in the raster setting. For a finite barrier raster $B$
with 8-connected barriers and a 4-connected complement, the target identity is
\[
  \#\{\text{enclosed 4-components of } \mathbb{Z}^2\setminus B\}
  \;=\; C_8(B) - \chi(K(B)).
\]
The proposed route is a ladder of three lemmas. L1 is a local formula for
$\chi(K(B))$. L2 says that adding one pixel preserves
$\text{enclosed} - C_8 + \chi$. L3 is the bridge itself, by induction on $|B|$.
The first version of this paper did L1: it gives a Lean proof of Gray's
bit-quad formula in its 8-connected form (Sections~\ref{sec:l1start}--\ref{sec:l1end}).
This version adds exact Lean statements, proposed for sealing, for L2 and L3
(Section~\ref{sec:proposals}). It also adds machine-checked sanity checks on
their definitions and computational evidence that the proposed statements are
true (Section~\ref{sec:evidence}). The 2026-11-04 revision made the proposal
modules self-contained (Section~\ref{sec:selfcontained}) and added an exhaustive
kernel check of all 256 neighbourhood configurations of L2
(Section~\ref{sec:allconfigs}). A proposed statement is a target, not a
result: those sections do not prove L2 or L3. L1 matters for the rest of the ladder because it makes the
change in $\chi$ under a single-pixel insertion a purely local quantity: only
the four windows around the new pixel change.

Both proposed statements were later sealed unchanged. The 2026-11-08 revision
(Section~\ref{sec:reduction}) reduced L2 to one direction of a half-plane digital
Jordan lemma about the pixel below the inserted one. It proved the other direction,
and with it half of the bridge identity, $\mathrm{enc}\ge C_8-\chi$. The present
revision closes the gap (Section~\ref{sec:crack}). We give a Lean proof of
the sealed statement \texttt{P1.Insertion} and, as a corollary, of the bridge
equation $\mathrm{enc}=C_8-\chi$ for every finite raster. The 2026-11-10 revision
(Section~\ref{sec:l3}) gives a Lean proof of the sealed L3 statement \texttt{P1.Bridge}
in its exact form. It also proves that exactly one component is unbounded, and states
precisely how far this goes towards the \texttt{FaceCount} hypothesis of
\emph{Trees Cannot Cut}.

\section{Background and related work}

Gray~\cite{Gray1971} introduced locally countable properties of binary images,
among them area, perimeter and the Euler number. He showed that the Euler number
can be computed from counts of $2\times2$ ``bit-quad'' patterns, with different
weightings for 4- and 8-connectivity of the foreground. The (8,4) pairing of
connectivity for foreground and background goes back to
Rosenfeld~\cite{Rosenfeld1970}, who showed that with this pairing a simple closed
digital curve separates the plane into exactly two components. Its role in
digital Jordan-type theorems is surveyed by Kong and
Rosenfeld~\cite{KongRosenfeld1989}.
In the 8-connected reading, two pixels that share only a corner touch in $K(B)$.
The two diagonal patterns then contribute $-2$ to $4\chi$, because their
shared corner is one vertex where four edges meet. Our Lean development uses
Lean~4~\cite{Lean4} and Mathlib~\cite{Mathlib2020}. We are not aware of an
earlier machine-checked proof of the bit-quad formula, but we have not searched
exhaustively and make no claim of novelty.

\section{Setting: the sealed statement}\label{sec:l1start}

The statement module \texttt{ViridisStatements.P1.BitQuad} fixes the following
definitions. A pixel $p=(i,j)\in\mathbb{Z}^2$ stands for the closed square
$[i,i+1]\times[j,j+1]$.
\begin{itemize}
\item $\mathrm{vertices}(B)=\bigcup_{p\in B}\mathrm{corners}(p)$, where
  $\mathrm{corners}(i,j)=\{(i,j),(i+1,j),(i,j+1),(i+1,j+1)\}$.
\item $\mathrm{hEdges}(B)=\bigcup_{(i,j)\in B}\{(i,j),(i,j+1)\}$ (horizontal
  edges named by their left end) and
  $\mathrm{vEdges}(B)=\bigcup_{(i,j)\in B}\{(i,j),(i+1,j)\}$ (vertical edges named
  by their lower end).
\item The Euler characteristic, computed in $\mathbb{Z}$:
  \[\chi(B) = |\mathrm{vertices}(B)| - (|\mathrm{hEdges}(B)|+|\mathrm{vEdges}(B)|) + |B|.\]
\item The window at a lattice point $v=(x,y)$ consists of four bits,
  $a=[(x-1,y-1)\in B]$, $b=[(x,y-1)\in B]$, $c=[(x-1,y)\in B]$ and
  $d=[(x,y)\in B]$, with $n(v)=a+b+c+d$.
\item $Q_1$, $Q_3$ and $Q_D$ count the vertices $v$ with $n(v)=1$, with
  $n(v)=3$, and with $\{a,d\}$ or $\{b,c\}$ as the only set bits.
\end{itemize}
Every window that meets $B$ is centred on a vertex of $K(B)$, so counting over
the vertices counts every such window exactly once.

\begin{theorem}[\texttt{P1.BitQuad}]\label{thm:main}
For every finite $B\subset\mathbb{Z}^2$,
$4\,\chi(K(B)) = Q_1(B) - Q_3(B) - 2\,Q_D(B)$.
\end{theorem}

\section{Proof architecture}

Write $\mathcal V=\mathrm{vertices}(B)$, and use $[P]\in\{0,1\}$ for the indicator
of a proposition $P$.

\begin{lemma}[Edge membership]\label{lem:edges}
$(i,j)\in\mathrm{hEdges}(B) \iff (i,j)\in B \lor (i,j-1)\in B$, and
$(i,j)\in\mathrm{vEdges}(B) \iff (i,j)\in B \lor (i-1,j)\in B$.
\end{lemma}

\begin{lemma}[Shift counting]\label{lem:shift}
Let $S,\mathcal V\subset\mathbb{Z}^2$ be finite and $s\in\mathbb{Z}^2$. If $x+s\in\mathcal V$
for every $x\in S$, then $\sum_{v\in\mathcal V}[v-s\in S]=|S|$.
\end{lemma}
\begin{proof}
The map $v\mapsto v-s$ is a bijection from $\{v\in\mathcal V: v-s\in S\}$ onto
$S$. Its inverse is $x\mapsto x+s$, which lands in $\mathcal V$ by hypothesis.
\end{proof}

Lemma~\ref{lem:shift} does all of the double counting. Every pixel $x\in B$
has its four corners $x$, $x+(1,0)$, $x+(0,1)$ and $x+(1,1)$ in $\mathcal V$.
Applying the lemma with the four shifts gives
\[
  4F = \sum_{v\in\mathcal V} n(v).
\]
By Lemma~\ref{lem:edges}, both ends $x$ and $x+(1,0)$ of a horizontal edge are
corners of the pixel above or below it, and likewise for vertical edges.
Write the arms of $v$ as left $[v-(1,0)\in\mathrm{hEdges}]$, right
$[v\in\mathrm{hEdges}]$, down $[v-(0,1)\in\mathrm{vEdges}]$ and up
$[v\in\mathrm{vEdges}]$. Applying the lemma once more for each arm gives
\[
  2E = \sum_{v\in\mathcal V}\mathrm{arms}(v), \qquad V=\sum_{v\in\mathcal V}1 .
\]
Hence $4\chi = \sum_{v\in\mathcal V}\bigl(4-2\,\mathrm{arms}(v)+n(v)\bigr)$.
Lemma~\ref{lem:edges} expresses each arm through the window bits. The left arm
is $a\lor c$, the right arm $b\lor d$, the down arm $a\lor b$ and the up arm
$c\lor d$. The summand therefore depends only on $(a,b,c,d)$, and
Table~\ref{tab:patterns} evaluates it. Every vertex is a corner of some pixel,
which must be one of the four window cells, so the empty pattern never occurs.
On the other 15 patterns the summand equals
$[n=1]-[n=3]-2[\text{diagonal}]$. Summing over $\mathcal V$ proves
Theorem~\ref{thm:main}.

\begin{table}[h]
\centering
\begin{tabular}{lcccc}
\toprule
Pattern class & \#patterns & $n$ & arms & $4-2\,\mathrm{arms}+n$\\
\midrule
empty (cannot occur) & 1 & 0 & 0 & 4\\
one pixel & 4 & 1 & 2 & $1$\\
adjacent pair & 4 & 2 & 3 & $0$\\
diagonal pair & 2 & 2 & 4 & $-2$\\
three pixels & 4 & 3 & 4 & $-1$\\
full & 1 & 4 & 4 & $0$\\
\bottomrule
\end{tabular}
\caption{The local summand on the 16 window patterns.}
\label{tab:patterns}
\end{table}

\section{The Lean development}

The candidate file \texttt{proofs/P1.BitQuad.lean} (145 lines) follows the
architecture above.
\begin{itemize}
\item \texttt{mem\_vertices}, \texttt{mem\_corners}, \texttt{mem\_hEdges} and
  \texttt{mem\_vEdges} unfold the membership predicates (Lemma~\ref{lem:edges}).
\item \texttt{card\_shift} proves Lemma~\ref{lem:shift} with
  \texttt{Finset.card\_bij}. Its corollary \texttt{sum\_shift} restates it as a
  sum of $\{0,1\}$-valued integer indicators via \texttt{Finset.sum\_boole}.
\item In \texttt{main}, the eight instances of \texttt{sum\_shift} (four for
  faces, two for each edge direction) and the three expressions of $Q_1$, $Q_3$
  and $Q_D$ as indicator sums reduce the goal to a single
  \texttt{Finset.sum\_congr}. The pointwise goal is closed by
  \texttt{by\_cases} on the four window memberships followed by \texttt{simp},
  which also uses the fact that the window at a vertex is non-empty.
\end{itemize}
The brief suggested evaluating the 16 patterns with \texttt{decide}. That
proved unnecessary: once the window memberships are case-split, \texttt{simp}
closes all 16 branches directly. The file uses no \texttt{decide} on large
objects and no \texttt{native\_decide}. The project checker
(\texttt{viridis-lean-check}) reports that the file compiles under Lean
\texttt{v4.28.0} with the pinned Mathlib, matches the sealed statement, depends
only on \texttt{propext}, \texttt{Classical.choice} and \texttt{Quot.sound}, and
replays through the kernel. Whether the claim is accepted is decided by the
project's independent verifier; see its certification page. The complete file
follows.

\begin{lstlisting}
import Mathlib
import ViridisStatements.P1.BitQuad

namespace Candidate

open ViridisStatements.P1.BitQuad

/-- Membership in `vertices`. -/
theorem mem_vertices (B : Finset Pixel) (v : ℤ × ℤ) :
    v ∈ vertices B ↔ ∃ p ∈ B, v ∈ corners p := by
  unfold vertices
  simp only [Finset.mem_biUnion]

theorem mem_corners (p v : ℤ × ℤ) :
    v ∈ corners p ↔ v = (p.1, p.2) ∨ v = (p.1 + 1, p.2) ∨ v = (p.1, p.2 + 1) ∨
      v = (p.1 + 1, p.2 + 1) := by
  unfold corners
  simp only [Finset.mem_insert, Finset.mem_singleton]

/-- A horizontal edge is present iff the pixel above or below it is in `B`. -/
theorem mem_hEdges (B : Finset Pixel) (h : ℤ × ℤ) :
    h ∈ hEdges B ↔ (h.1, h.2) ∈ B ∨ (h.1, h.2 - 1) ∈ B := by
  unfold hEdges
  simp only [Finset.mem_biUnion, Finset.mem_insert, Finset.mem_singleton]
  constructor
  · rintro ⟨p, hp, rfl | rfl⟩
    · left; simpa using hp
    · right; simpa using hp
  · rintro (h1 | h1)
    · exact ⟨_, h1, Or.inl rfl⟩
    · exact ⟨_, h1, Or.inr (by simp)⟩

/-- A vertical edge is present iff the pixel right or left of it is in `B`. -/
theorem mem_vEdges (B : Finset Pixel) (e : ℤ × ℤ) :
    e ∈ vEdges B ↔ (e.1, e.2) ∈ B ∨ (e.1 - 1, e.2) ∈ B := by
  unfold vEdges
  simp only [Finset.mem_biUnion, Finset.mem_insert, Finset.mem_singleton]
  constructor
  · rintro ⟨p, hp, rfl | rfl⟩
    · left; simpa using hp
    · right; simpa using hp
  · rintro (h1 | h1)
    · exact ⟨_, h1, Or.inl rfl⟩
    · exact ⟨_, h1, Or.inr (by simp)⟩

/-- Shift-counting: if every translate `x + (a, b)` of `S` lies in `V`, then
counting the `v ∈ V` with `v - (a, b) ∈ S` gives `|S|`. -/
theorem card_shift (V S : Finset (ℤ × ℤ)) (a b : ℤ)
    (hS : ∀ x ∈ S, (x.1 + a, x.2 + b) ∈ V) :
    (V.filter fun v => (v.1 - a, v.2 - b) ∈ S).card = S.card := by
  apply Finset.card_bij (fun v _ => (v.1 - a, v.2 - b))
  · intro v hv
    exact (Finset.mem_filter.mp hv).2
  · intro v₁ _ v₂ _ h
    simp only [Prod.mk.injEq] at h
    ext <;> omega
  · intro x hx
    refine ⟨(x.1 + a, x.2 + b), Finset.mem_filter.mpr ⟨hS x hx, ?_⟩, ?_⟩
    · simpa using hx
    · simp

theorem sum_shift (V S : Finset (ℤ × ℤ)) (a b : ℤ)
    (hS : ∀ x ∈ S, (x.1 + a, x.2 + b) ∈ V) :
    (∑ v ∈ V, if (v.1 - a, v.2 - b) ∈ S then (1 : ℤ) else 0) = S.card := by
  rw [Finset.sum_boole, card_shift V S a b hS]

theorem main : ViridisStatements.P1.BitQuad.Statement := by
  intro B
  set V := vertices B with hV
  -- inclusion facts
  have hB : ∀ x ∈ B, ∀ a b : ℤ, (a = 0 ∨ a = 1) → (b = 0 ∨ b = 1) →
      (x.1 + a, x.2 + b) ∈ V := by
    intro x hx a b ha hb
    rw [hV, mem_vertices]
    refine ⟨x, hx, ?_⟩
    rw [mem_corners]
    rcases ha with rfl | rfl <;> rcases hb with rfl | rfl <;> simp
  have hH : ∀ x ∈ hEdges B, ∀ a : ℤ, (a = 0 ∨ a = 1) → (x.1 + a, x.2 + 0) ∈ V := by
    intro x hx a ha
    rw [mem_hEdges] at hx
    rcases hx with hx | hx
    · have := hB _ hx a 0 ha (Or.inl rfl)
      simpa using this
    · have := hB _ hx a 1 ha (Or.inr rfl)
      simpa using this
  have hVe : ∀ x ∈ vEdges B, ∀ b : ℤ, (b = 0 ∨ b = 1) → (x.1 + 0, x.2 + b) ∈ V := by
    intro x hx b hb
    rw [mem_vEdges] at hx
    rcases hx with hx | hx
    · have := hB _ hx 0 b (Or.inl rfl) hb
      simpa using this
    · have := hB _ hx 1 b (Or.inr rfl) hb
      simpa using this
  -- the four expressions as sums over V
  have eV : ((V.card : ℤ)) = ∑ v ∈ V, (1 : ℤ) := by simp
  have eF11 := sum_shift V B 1 1 (fun x hx => hB x hx 1 1 (Or.inr rfl) (Or.inr rfl))
  have eF01 := sum_shift V B 0 1 (fun x hx => hB x hx 0 1 (Or.inl rfl) (Or.inr rfl))
  have eF10 := sum_shift V B 1 0 (fun x hx => hB x hx 1 0 (Or.inr rfl) (Or.inl rfl))
  have eF00 := sum_shift V B 0 0 (fun x hx => hB x hx 0 0 (Or.inl rfl) (Or.inl rfl))
  have eH10 := sum_shift V (hEdges B) 1 0 (fun x hx => hH x hx 1 (Or.inr rfl))
  have eH00 := sum_shift V (hEdges B) 0 0 (fun x hx => hH x hx 0 (Or.inl rfl))
  have eV01 := sum_shift V (vEdges B) 0 1 (fun x hx => hVe x hx 1 (Or.inr rfl))
  have eV00 := sum_shift V (vEdges B) 0 0 (fun x hx => hVe x hx 0 (Or.inl rfl))
  have eQ1 : ((Q1 B : ℕ) : ℤ) = ∑ v ∈ V, if windowCount B v = 1 then (1 : ℤ) else 0 := by
    rw [Finset.sum_boole]; rfl
  have eQ3 : ((Q3 B : ℕ) : ℤ) = ∑ v ∈ V, if windowCount B v = 3 then (1 : ℤ) else 0 := by
    rw [Finset.sum_boole]; rfl
  have eQD : ((QD B : ℕ) : ℤ) = ∑ v ∈ V, if isDiagonal B v = true then (1 : ℤ) else 0 := by
    rw [Finset.sum_boole]; rfl
  have key : 4 * euler B = ∑ v ∈ V,
      ((4 : ℤ) - 2 * ((if (v.1 - 1, v.2 - 0) ∈ hEdges B then 1 else 0)
        + (if (v.1 - 0, v.2 - 0) ∈ hEdges B then 1 else 0)
        + (if (v.1 - 0, v.2 - 1) ∈ vEdges B then 1 else 0)
        + (if (v.1 - 0, v.2 - 0) ∈ vEdges B then 1 else 0))
        + ((if (v.1 - 1, v.2 - 1) ∈ B then 1 else 0)
        + (if (v.1 - 0, v.2 - 1) ∈ B then 1 else 0)
        + (if (v.1 - 1, v.2 - 0) ∈ B then 1 else 0)
        + (if (v.1 - 0, v.2 - 0) ∈ B then 1 else 0))) := by
    simp only [Finset.sum_add_distrib, Finset.sum_sub_distrib, ← Finset.mul_sum]
    rw [eF11, eF01, eF10, eF00, eH10, eH00, eV01, eV00]
    unfold euler
    rw [← hV]
    simp
    ring
  rw [key, eQ1, eQ3, eQD, ← Finset.sum_sub_distrib, Finset.mul_sum,
    ← Finset.sum_sub_distrib]
  apply Finset.sum_congr rfl
  intro v hv
  obtain ⟨x, y⟩ := v
  have hne : (x - 1, y - 1) ∈ B ∨ (x, y - 1) ∈ B ∨ (x - 1, y) ∈ B ∨ (x, y) ∈ B := by
    rw [hV, mem_vertices] at hv
    obtain ⟨p, hp, hc⟩ := hv
    rw [mem_corners] at hc
    rcases hc with h | h | h | h <;> simp only [Prod.mk.injEq] at h <;> obtain ⟨rfl, rfl⟩ := h
    · right; right; right; exact hp
    · right; right; left; simpa using hp
    · right; left; simpa using hp
    · left; simpa using hp
  unfold windowCount isDiagonal bits
  simp only [sub_zero, mem_hEdges, mem_vEdges]
  by_cases ha : (x - 1, y - 1) ∈ B <;> by_cases hb : (x, y - 1) ∈ B <;>
    by_cases hc : (x - 1, y) ∈ B <;> by_cases hd : (x, y) ∈ B <;>
    simp [ha, hb, hc, hd] at hne ⊢

end Candidate
\end{lstlisting}

\section{Results}

\begin{itemize}
\item We give a Lean proof of Theorem~\ref{thm:main}, the sealed statement
  \texttt{P1.BitQuad}, for all finite rasters.
\item As an independent computational check, which is evidence and not a
  proof, the script \texttt{code/bitquad\_check.py} recomputes both sides for
  all $2^{16}=65{,}536$ subsets of a $4\times4$ box. It finds no
  counterexample (\texttt{results/bitquad\_4x4.json}). Euler characteristics
  from $-1$ to $4$ occur in that box.
\end{itemize}

\section{How L1 serves the L2/L3 bridge}\label{sec:l1end}

Let $p\notin B$ and $B'=B\cup\{p\}$. The windows whose bits differ between $B$
and $B'$ are exactly the four windows centred on the corners of $p$. By
Theorem~\ref{thm:main} applied to $B$ and to $B'$, $4(\chi(B')-\chi(B))$ equals
the change in $Q_1-Q_3-2Q_D$ over those four windows. That change depends only on the
$3\times3$ neighbourhood of $p$. This is the ``local half'' of L2. The non-local
half, the change in $C_8$ and in the number of enclosed 4-components of the
complement, still needs a separate argument: a pixel joining two
8-components can close off a new bounded region. L1 contributes nothing there.
The negative controls listed in the project context (the solid $2\times2$ block,
whose 8-neighbour graph $K_4$ has circuit rank 3 yet encloses nothing, and the
four-pixel N/S/E/W cross enclosing its centre) show that L2 must be stated
in terms of $\chi(K(B))$ and not in terms of the graph circuit rank.


\section{Proposed statements for L2 and L3}\label{sec:proposals}

The two modules \texttt{proposals/P1.Insertion.lean} and
\texttt{proposals/P1.Bridge.lean} are \emph{proposals}. They become targets only
if they are sealed, and nothing in this section is a proof of either statement.
Proposals may not import one another, so the two modules carry identical copies
of the connectivity definitions.

\subsection{Self-contained modules}\label{sec:selfcontained}
{\sloppy The first draft of these modules (2026-11-03) imported the sealed L1 module
\texttt{ViridisStatements.P1.BitQuad} for $\chi$. In the verifier's proposal
environment that import did not resolve: the compiled object file for the module
was missing. Both drafts were therefore rejected before any probe ran. The brief
anticipates this: if the L1 module cannot be imported, its definitions are
restated identically. Both modules now import only Mathlib. Each restates
\texttt{Pixel}, \texttt{corners}, \texttt{vertices}, \texttt{hEdges},
\texttt{vEdges} and \texttt{euler} character for character from the sealed L1
module, and a textual \texttt{diff} confirms that the copies are identical. The
two \texttt{Statement} bodies are unchanged from the first draft. If the import
becomes available, the restated block can be replaced by
\texttt{open ViridisStatements.P1.BitQuad} without any change in meaning.\par}

\subsection{Definitions and conventions}

\begin{definition}
Let $G_4$ and $G_8$ be the graphs on $\mathbb{Z}^2$ in which distinct $p,q$ are
adjacent when $|p_1-q_1|+|p_2-q_2|=1$, respectively when $|p_1-q_1|\le1$ and
$|p_2-q_2|\le1$. For a finite $B$, $C_8(B)$ is the number of connected
components of the subgraph of $G_8$ induced on $B$. A connected component of the
subgraph of $G_4$ induced on $\mathbb{Z}^2\setminus B$ is \emph{enclosed} when it
is finite, and $\mathrm{enc}(B)$ is the number of enclosed components.
\end{definition}

The brief for this task fixes the exterior convention, and the definitions follow
it literally. The complement is taken in all of $\mathbb{Z}^2$, never inside a
cropping window. ``Enclosed'' means ``finite'', so no bounding box enters the
statements. In Lean the graphs are built with \texttt{SimpleGraph.fromRel},
which symmetrises a relation and removes loops, and components are Mathlib's
\texttt{ConnectedComponent} of an induced graph
(Listing~\ref{lst:defs}).

\begin{lstlisting}[caption={Definitions and the two proposed statements (each Statement lives in its own module).},label={lst:defs}]
def G4 : SimpleGraph Pixel :=
  SimpleGraph.fromRel fun p q => |p.1 - q.1| + |p.2 - q.2| = 1
def G8 : SimpleGraph Pixel :=
  SimpleGraph.fromRel fun p q => |p.1 - q.1| ≤ 1 ∧ |p.2 - q.2| ≤ 1
def barrierGraph (B : Finset Pixel) : SimpleGraph (B : Set Pixel) :=
  G8.induce (B : Set Pixel)
def complementGraph (B : Finset Pixel) : SimpleGraph ((B : Set Pixel)ᶜ : Set Pixel) :=
  G4.induce (B : Set Pixel)ᶜ
noncomputable def C8 (B : Finset Pixel) : ℕ :=
  Nat.card (barrierGraph B).ConnectedComponent
def IsEnclosed {B : Finset Pixel} (c : (complementGraph B).ConnectedComponent) : Prop :=
  c.supp.Finite
noncomputable def enclosed (B : Finset Pixel) : ℕ :=
  Nat.card {c : (complementGraph B).ConnectedComponent // IsEnclosed c}

-- P1.Bridge (L3)
def Statement : Prop :=
  ∀ B : Finset Pixel,
    {c : (complementGraph B).ConnectedComponent | IsEnclosed c}.Finite ∧
      (enclosed B : ℤ) = C8 B - euler B

-- P1.Insertion (L2)
noncomputable def invariant (B : Finset Pixel) : ℤ :=
  (enclosed B : ℤ) - C8 B + euler B
def Statement : Prop :=
  ∀ (B : Finset Pixel) (p : Pixel), p ∉ B → invariant (insert p B) = invariant B
\end{lstlisting}

\paragraph{L3 (\texttt{P1.Bridge}).} For every finite $B$, the set of enclosed
components is finite and $\mathrm{enc}(B) = C_8(B)-\chi(K(B))$ in $\mathbb{Z}$.
The finiteness clause is deliberate. \texttt{Nat.card} returns $0$ on an
infinite type, so without the clause the equation could be met in a degenerate
way. With it, the count means what it says. The clause is true because every
component outside the bounding box of $B$ contains a ray.

\paragraph{L2 (\texttt{P1.Insertion}).} For every finite $B$ and every
$p\notin B$, $\mathrm{enc}-C_8+\chi$ takes the same value on $B\cup\{p\}$ and on
$B$. The statement puts no condition on the $3\times3$ neighbourhood of $p$. All
256 configurations of its eight neighbours are therefore covered, in every global
context. These include insertions that merge several 8-components, insertions
that close off a hole, and insertions that fill one. By L1, $\Delta\chi$ is a
function of the neighbourhood (Section~\ref{sec:l1end}). $\Delta C_8$ and
$\Delta\mathrm{enc}$ are not local, and they are where the proof effort lies.

\paragraph{Relation between the rungs.} L3's equation follows from L2 by
induction on $|B|$ from the empty raster, where both sides are $0$. Conversely,
L3's equation gives L2 at once. The two rungs are equivalent given the base case.
The split organises the proof around one inductive step; it does not grade the
statements by strength.

\subsection{Controls}

The definitions of $C_8$ and $\mathrm{enc}$ are not decidable, since they
quantify over infinite components of $\mathbb{Z}^2$. Each module therefore also
contains a computable \emph{bounded-box model}. It flood-fills the complement
inside the bounding box of $B\cup\{(0,0)\}$ with a one-pixel margin, starting
from the box corner, and counts components by their lexicographically least
pixel. The statements do not mention this model, and we do not prove that it
agrees with the definitions. Against the box model, \texttt{decide} checks the
values in Table~\ref{tab:controls}. The \texttt{P1.Insertion} module also checks
four non-local insertions: closing the ring's top edge, merging the ring's two
columns through a corner, completing the diamond, and filling the ring's centre.

\begin{table}[h]
\centering
\begin{tabular}{lcccl}
\toprule
raster & $C_8$ & $\chi$ & enc & role\\
\midrule
empty & 0 & 0 & 0 & base case\\
$3\times3$ ring & 1 & 0 & 1 & witness\\
solid $2\times2$ block & 1 & 1 & 0 & negative control 1 ($K_4$ has circuit rank 3)\\
N/S/E/W diamond & 1 & 0 & 1 & negative control 2 (4-adjacency graph has no edges)\\
diagonal pair & 1 & 1 & 0 & corner contact\\
\bottomrule
\end{tabular}
\caption{Control values. All rows are checked by \texttt{decide} on the box
model. For the first four rows we also give Lean proofs for the semantic
definitions $C_8$ and $\mathrm{enc}$ themselves.}
\label{tab:controls}
\end{table}

To connect the controls to the definitions that the statements actually use, each
module proves the control values for $C_8$ and $\mathrm{enc}$ directly. Two
general lemmas carry these proofs (Listing~\ref{lst:controls}). If a horizontal
ray from a non-barrier pixel avoids $B$, the pixel's component is infinite
(\texttt{infinite\_of\_ray}). A non-barrier pixel whose four 4-neighbours lie in
$B$ forms a finite component (\texttt{isolated}). For the ring, every non-barrier
pixel other than the centre has a free ray to the left or right. The centre is
isolated. Hence exactly one component is enclosed. The block and diamond are
handled the same way, and $C_8=1$ follows from explicit 8-paths. Combined with
\texttt{decide} on $\chi$, this checks the identity $\mathrm{enc}=C_8-\chi$ for
the ring ($1=1-0$), the block ($0=1-1$) and the diamond ($1=1-0$) on the
semantic definitions.

\begin{lstlisting}[caption={Key lemmas behind the semantic controls (excerpt from the proposal modules).},label={lst:controls}]
/-- A horizontal ray of non-barrier pixels makes a component infinite. -/
theorem infinite_of_ray (B : Finset Pixel) (v : Pixel) (hv : v ∉ B) (d : ℤ) (hd : d = 1 ∨ d = -1)
    (hray : ∀ n : ℕ, (v.1 + d * n, v.2) ∉ B) :
    ¬ IsEnclosed ((complementGraph B).connectedComponentMk ⟨v, hv⟩) := by
  intro hfin
  have hreach : ∀ n : ℕ, (complementGraph B).Reachable ⟨v, hv⟩ ⟨(v.1 + d * n, v.2), hray n⟩ := by
    intro n
    induction n with
    | zero => simp
    | succ n ih =>
      refine ih.trans (SimpleGraph.Adj.reachable ?_)
      simp only [complementGraph, SimpleGraph.comap_adj, SimpleGraph.induce, G4,
        SimpleGraph.fromRel_adj, Function.Embedding.coe_subtype, ne_eq, Prod.mk.injEq]
      push_cast
      rcases hd with rfl | rfl <;> refine ⟨by omega, Or.inl ?_⟩ <;> simp
  apply Set.infinite_of_injective_forall_mem (f := fun n : ℕ => (⟨(v.1 + d * n, v.2), hray n⟩ :
    ((B : Set Pixel)ᶜ : Set Pixel)))
  · intro m n h
    simp only [Subtype.mk.injEq, Prod.mk.injEq] at h
    rcases hd with rfl | rfl <;> omega
  · intro n
    exact (SimpleGraph.ConnectedComponent.mem_supp_iff _ _).2
      (SimpleGraph.ConnectedComponent.eq.2 (hreach n).symm)
  · exact hfin


/-- A pixel whose four 4-neighbours are all barrier is a finite component on its own. -/
theorem isolated (B : Finset Pixel) (h : Pixel) (hh : h ∉ B)
    (hn : ∀ u : Pixel, |h.1 - u.1| + |h.2 - u.2| = 1 → u ∈ B) :
    IsEnclosed ((complementGraph B).connectedComponentMk ⟨h, hh⟩) := by
  apply (Set.finite_singleton (⟨h, hh⟩ : ((B : Set Pixel)ᶜ : Set Pixel))).subset
  intro w hw
  rw [SimpleGraph.ConnectedComponent.mem_supp_iff, SimpleGraph.ConnectedComponent.eq] at hw
  obtain ⟨p⟩ := hw.symm
  cases p with
  | nil => rfl
  | @cons _ u _ hadj _ =>
    exfalso
    simp only [complementGraph, SimpleGraph.comap_adj, SimpleGraph.induce, G4,
      SimpleGraph.fromRel_adj, Function.Embedding.coe_subtype] at hadj
    obtain ⟨-, hr | hr⟩ := hadj
    · exact u.2 (hn _ hr)
    · exact u.2 (hn _ (by rw [abs_sub_comm h.1, abs_sub_comm h.2]; exact hr))

theorem enclosed_ring : enclosed ring = 1 := by
  refine enclosed_eq_one ring (1, 1) (by decide) (isolated _ _ _ ?_) ?_
  · rintro ⟨x, y⟩ h
    have h' : |1 - x| + |1 - y| = 1 := h
    simp only [ring, Finset.mem_insert, Finset.mem_singleton, Prod.mk.injEq]
    rcases abs_cases (1 - x) with ⟨h1, h1'⟩ | ⟨h1, h1'⟩ <;>
    rcases abs_cases (1 - y) with ⟨h2, h2'⟩ | ⟨h2, h2'⟩ <;> omega
  · rintro ⟨x, y⟩ hv hne
    simp only [ring, Finset.mem_insert, Finset.mem_singleton, Prod.mk.injEq, not_or] at hv
    simp only [ne_eq, Prod.mk.injEq] at hne
    by_cases hr : 2 < x ∨ y < 0 ∨ 2 < y
    · refine infinite_of_ray _ _ _ 1 (Or.inl rfl) (ray_right _ _ ?_)
      simp only [ring, Finset.mem_insert, Finset.mem_singleton, forall_eq_or_imp, forall_eq]
      omega
    · refine infinite_of_ray _ _ _ (-1) (Or.inr rfl) (ray_left _ _ ?_)
      simp only [ring, Finset.mem_insert, Finset.mem_singleton, forall_eq_or_imp, forall_eq]
      omega

\end{lstlisting}

\subsection{All 256 neighbourhood configurations by \texttt{decide}}\label{sec:allconfigs}

The task brief asks that L2 account for every configuration of the $3\times3$
neighbourhood. The statement itself does so by having no hypothesis on the
neighbourhood. The first draft checked only four hand-picked insertions in Lean,
because the \texttt{Finset} flood fill of the box model is too slow for the kernel
to run through all 256 cases. This revision adds a second test oracle (namespace
\texttt{Mask} in both modules). A $7\times7$ box is a 49-bit natural number with
pixel $(x,y)$ at bit $7y+x$. Neighbourhoods are computed with shifts and masks,
components by iterated flood fill from the lowest set bit, and $V-E+F$ by
counting the bits of shifted unions. These use the same vertex and edge naming as
\texttt{euler}. The Lean kernel evaluates these natural-number operations
efficiently, so \texttt{decide} handles the full table
(Listing~\ref{lst:mask}). With the inserted pixel $p$ at the centre $(3,3)$:
\begin{itemize}
\item \texttt{allConfigs 0 256}: for each of the 256 subsets $N$ of the eight
  neighbours of $p$, the invariant is $0$ on $N$ and on $N\cup\{p\}$. So all 512
  subsets of the $3\times3$ window satisfy the L3 identity, and inserting the
  centre never changes $\mathrm{enc}-C_8+\chi$.
\item \texttt{allConfigs ring5 256}: the same for $N$ placed inside the closed
  16-pixel ring at Chebyshev distance 2 from $p$. Every barrier neighbour of $p$
  is then already 8-connected to every other through the ring, and every
  non-barrier run around $p$ lies in a finite region. This is the opposite
  extreme of global context from the first check.
\end{itemize}
The bitmask model also reproduces the witness and both negative controls,
including the four 4-components of the diamond. In a scratch mutation test (not
part of the modules), we replaced 8-adjacency by 4-adjacency in the model's
$C_8$. Both exhaustive checks then failed, so the table depends on the (8,4)
pairing. As with the box model, we do not prove that the bitmask model agrees
with the semantic definitions. Two contexts are also far from all contexts: the
proof of L2 must still handle every intermediate global situation.

\begin{lstlisting}[caption={The exhaustive configuration check (excerpt from the \texttt{Mask} namespace).},label={lst:mask}]
def inv (B : ℕ) : ℤ := (enc B : ℤ) - c8 B + chi B
def allConfigs (ctx : ℕ) : ℕ → Bool
  | 0 => true
  | n + 1 =>
    let B := ctx ||| spread n
    (inv (B ||| centre) == inv B && inv B == 0) && allConfigs ctx n

set_option maxRecDepth 100000 in
example : allConfigs 0 256 = true := by decide
set_option maxRecDepth 100000 in
example : allConfigs ring5 256 = true := by decide
\end{lstlisting}

\section{A proof of L2, part 1: reduction to a half-plane Jordan lemma}\label{sec:reduction}

This section reports the 2026-11-08 session on the sealed statement
\texttt{P1.Insertion}. Section~\ref{sec:crack} completes the proof. The candidate file
\texttt{proofs/P1.Insertion.lean} (about 1{,}950 lines including the 2026-11-09
additions) compiles against the pinned toolchain, and its theorems depend only on
\texttt{propext}, \texttt{Classical.choice} and \texttt{Quot.sound}. Write $B' = B\cup\{p\}$ with $p\notin B$, and
$W_p(B) = \{q\in B : |q_1-p_1|\le1,\ |q_2-p_2|\le1\}$ for the part of $B$ in the
$3\times3$ window of $p$.

\paragraph{Locality of $\Delta\chi$.} We give a Lean proof (\texttt{euler\_insert\_local}) that
\[
  \chi(K(B')) - \chi(K(B)) \;=\; \chi(K(W_p(B)\cup\{p\})) - \chi(K(W_p(B))).
\]
The proof does not go through L1. Each of \texttt{vertices}, \texttt{hEdges} and
\texttt{vEdges} is a union over pixels, so inserting $p$ adds exactly the cells of
$p$ that are not already present. A corner or edge of $p$ can belong to a pixel $q$
only if $q$ is in the window, so whether such a cell is new is decided by $W_p(B)$.

\paragraph{The change in $C_8$.} Let $k$ be the number of 8-components of $B$ that
contain an 8-neighbour of $p$. We give a Lean proof (\texttt{C8\_insert}) that
$C_8(B') = C_8(B) + 1 - k$. The inclusion $B\subseteq B'$ is a graph homomorphism of the
8-adjacency graphs, so it maps components to components. The key lemma is a walk
lemma, proved by induction on \texttt{Relation.ReflTransGen}. If $v\in B$ reaches
$y$ in $B'$, then either $y\ne p$ and $v$ reaches $y$ already in $B$, or the component
of $v$ touches $p$. From this, the components that touch $p$ all map to the
component of $p$, and the others map injectively onto the rest. The bijection
$\mathrm{Option}\{\text{non-touching}\}\simeq \pi_0(B')$ gives the count.

\paragraph{The change in the enclosed count.} The complement of $B'$ is the
complement of $B$ with $p$ removed, and the same walk lemma applies with
4-adjacency. Let $c_p$ be the 4-component of $p$ in $\mathbb{Z}^2\setminus B$ and
let $P$ be the number of enclosed 4-components of $\mathbb{Z}^2\setminus B'$ that
contain a 4-neighbour of $p$. A component of $\mathbb{Z}^2\setminus B'$ that
contains no 4-neighbour of $p$ has the same support as its image, so it is enclosed exactly
when its image is. We give a Lean proof (\texttt{enclosed\_insert}) that, if $B$ has finitely many enclosed
components,
\[
  \mathrm{enc}(B') = \mathrm{enc}(B) - [\,c_p \text{ enclosed}\,] + P .
\]
The finiteness hypothesis is needed because \texttt{enclosed} is a
\texttt{Nat.card}, which is $0$ on infinite types.

\paragraph{Finiteness.} We give a Lean proof (\texttt{finite\_enclosed}) that for every finite $B$
the set of enclosed components of $\mathbb{Z}^2\setminus B$ is finite. Choose
$M$ with $|b_1|,|b_2|<M$ for all $b\in B$. A non-barrier pixel $v$ with
$v_1\ge M$, $v_1\le -M$ or $|v_2|\ge M$ has a horizontal ray free of $B$, so by the
lemma \texttt{infinite\_of\_ray} of the sealed module its component is infinite.
Hence every enclosed component has a representative in the finite box
$[-M,M]^2$, and sending a component to its representative is injective. This also
discharges the finiteness hypothesis above, and it is clause (1) of the sealed
statement \texttt{P1.Bridge}. It is stated there in a different form, so we do
not claim that statement.

\paragraph{The reduction.} Combining the four gives \texttt{invariant\_insert}:
\[
  \Delta(\mathrm{enc}-C_8+\chi) \;=\; \bigl(P - [\,c_p\text{ enclosed}\,]\bigr) - (1-k)
  + \Delta\chi\bigl(W_p(B)\bigr),
\]
and \texttt{statement\_of\_jordan} derives \texttt{Statement} from one hypothesis:
\begin{enumerate}
\item[(J)] for every $B$ and $p\notin B$,
  $P - [\,c_p\text{ enclosed}\,] - (1-k) + \Delta\chi(W_p(B)) = 0$.
\end{enumerate}
(J) is the real topological content. Let $a$ be the number of
maximal 8-connected arcs of barrier pixels on the 8-cycle around $p$. For
$r\ge1$ non-barrier 4-neighbours we expect $\Delta\chi(W_p)=1-a$, and $P$ should
equal $a-k+1$ minus the unbounded piece when $c_p$ is infinite. Each pair of arcs that is
already joined by a barrier 8-path away from $p$ closes a digital curve through $p$
and cuts off one more piece. Proving this needs a digital Jordan curve theorem for
(8,4) curves~\cite{Rosenfeld1970}. Mathlib does not provide one. The rest of this section
narrows (J) down to a single configuration.

\paragraph{A special case of L2.} As a check on the reduction, we give a Lean proof
(\texttt{insertion\_hole}) that $\mathrm{enc}-C_8+\chi$ is unchanged when all four
4-neighbours of $p$ are barrier pixels. In this case no new component contains a
4-neighbour of $p$, so $P=0$. The old component $c_p=\{p\}$ is enclosed, by the
lemma \texttt{isolated} of the sealed module. Every 8-neighbour of $p$ is joined to
the north neighbour through the north, east and south neighbours, so $k=1$. Every
corner and edge of $p$ is already present, so $\Delta\chi=1$. Then (J) reads
$0-1-0+1=0$. This case covers, for example, filling the centre of the $3\times3$
ring and of the diamond. It is not the hard case: there, merging and splitting
interact through the global shape of $B$.


\paragraph{Raster-scan order.} Say that $B$ is \emph{before} $p$ if every
$q\in B$ has $q_2<p_2$, or $q_2=p_2$ and $q_1<p_1$. The invariant of the empty
raster is $0$, by the control lemmas of the sealed module and $\chi(\emptyset)=0$.
Mathlib's \texttt{Finset.induction\_on\_max\_value} with the lexicographic key
$(q_2,q_1)$ builds every finite raster by insertions in which $B$ is before $p$. So
if those insertions preserve the invariant, the invariant is $0$ on every finite
raster, and \texttt{Statement} follows (\texttt{statement\_of\_scan}). This is the
classical raster-scan idea of connected-component labelling~\cite{RosenfeldPfaltz1966}.

\paragraph{Insertion in scan order.} Write $W,SW,S,SE$ for $p+(-1,0)$, $p+(-1,-1)$,
$p+(0,-1)$, $p+(1,-1)$. If $B$ is before $p$, every pixel above the row of $p$, and
every pixel to the right of $p$ in its row, is free in $B\cup\{p\}$ and has a free
ray. We give Lean proofs of the following. The old component $c_p$ is infinite. A
new piece that contains $E$, $N$ or $W$ is infinite: $W$ is 4-adjacent to the free
pixel $NW$. So $P$ is $1$ if $S$ is free and its piece is enclosed, and $0$ otherwise
(\texttt{P\_scan}). The barrier neighbours of $p$ lie among $W,SW,S,SE$, and
(\texttt{euler\_scan})
\[
  \Delta\chi = [\,W,S,SW\notin B\,] + [\,S,SE\notin B\,] - [\,S\notin B\,].
\]
Table~\ref{tab:scan} lists the five cases. In four of them $k$ and $P$ are fixed
by local arguments, and the invariant is preserved. In the fifth, $S$ is free,
$SE\in B$ and $L\in\{W,SW\}$ meets $B$. Then $\Delta\chi=-1$, $k\in\{1,2\}$, and we
give a Lean proof (\texttt{k\_core}) that $k=1$ exactly when $L$ and $SE$ are
8-connected in $B$.

\begin{table}[h]
\centering
\begin{tabular}{@{}lcccl@{}}
\toprule
Case (scan order) & $\Delta\chi$ & $k$ & $P$ & reason for $P$\\
\midrule
$S\in B$ & $0$ & $1$ & $0$ & $S$ is not free\\
$S,SE\notin B$, $W,SW\notin B$ & $1$ & $0$ & $0$ & $S\to SE\to E$, free ray\\
$S,SE\notin B$, $L\cap B\neq\emptyset$ & $0$ & $1$ & $0$ & $S\to SE\to E$\\
$S\notin B$, $SE\in B$, $W,SW\notin B$ & $0$ & $1$ & $0$ & $S\to SW\to W\to NW$\\
$S\notin B$, $SE\in B$, $L\cap B\neq\emptyset$ & $-1$ & $1$ or $2$ & $2-k$ & half-plane Jordan lemma\\
\bottomrule
\end{tabular}
\caption{Insertion of $p$ after all of $B$ in raster-scan order. In every row the
invariant changes by $P-(1-k)+\Delta\chi$ (the old component of $p$ is never
enclosed), so it is preserved exactly when $P+k+\Delta\chi=1$. Rows 1--4 were proved
by local arguments on 2026-11-08. Row 5 needs $P+k=2$, which is the half-plane Jordan
lemma (crossing parity, Section~\ref{sec:reduction}, and the crack boundary,
Section~\ref{sec:crack}).}
\label{tab:scan}
\end{table}


\paragraph{One direction by crossing parity.} We give a Lean proof
(\texttt{senc\_of\_join}) of the direction ``$L$ and $SE$ 8-connected
$\Rightarrow$ the component of $S$ is finite''. Take an 8-walk $\gamma$ in $B$ from
$L$ to $SE$ and close it into the cycle $\Gamma=(p,\gamma,p)$. For a pixel $v$ let
$N(v)$ count the steps $a\to b$ of $\Gamma$ that join column $v_1$ to column
$v_1+1$ and whose endpoint in column $v_1$ lies strictly above $v$. Each of the
following facts reduces to an identity, modulo $2$, for a single step of $\Gamma$.
Each such identity is a statement in linear integer arithmetic, closed by
\texttt{omega}. A generic lemma (\texttt{dsum\_parity}) then adds them up along the
walk:
\begin{itemize}
\item if $v$ and $v+(0,1)$ are free, $N(v)\equiv N(v+(0,1))$: a step could
  distinguish them only through the free pixel $v+(0,1)$;
\item if $v$ and $v+(1,0)$ are free, $N(v)+N(v+(1,0))$ is the number of times
  $\Gamma$ enters or leaves the half-column $\{(v_1+1,t):t>v_2\}$, which is even
  because $\Gamma$ is closed;
\item $N(v)$ is even if $v$ is above the row of $p$, to the left or right of all of
  $\Gamma$, or below all of $\Gamma$ (then $N(v)$ counts every crossing of a full vertical
  line, which is even);
\item $N(S)$ is odd: the only step of $\Gamma$ that reaches column $x+1$ from a pixel
  of column $x$ at height at least $y$ is $SE\to p$, because $B$ is before $p$.
\end{itemize}
So $N$ is odd on the whole free 4-component of $S$, which therefore lies in a
finite box.

\paragraph{Half of the bridge identity.} In the fifth case the invariant changes by
$P+k-2$. If $L$ and $SE$ are connected then $k=1$ and, by \texttt{senc\_of\_join},
$P=1$. Otherwise $k=2$ and $P\ge0$. So every scan-order insertion changes the
invariant by $0$ or $+1$ (\texttt{scan\_step\_mono}). Since the empty raster has
invariant $0$, we give a Lean proof (\texttt{invariant\_nonneg}, \texttt{enclosed\_ge})
that for every finite raster
\[
  \mathrm{enc}(B) \;\ge\; C_8(B) - \chi(K(B)).
\]
The invariant of $B$ counts the ``bad'' scan steps, where $S$ is enclosed although $L$
and $SE$ lie in different 8-components. \texttt{statement\_of\_le} derives
\texttt{Statement} from the reverse inequality.

\paragraph{What remained on 2026-11-08.} \texttt{statement\_of\_halfplane\_jordan} derives
\texttt{Statement} from the following hypothesis. From the
same hypothesis, \texttt{bridge\_of\_halfplane\_jordan} derives the equation
$\mathrm{enc}(B)=C_8(B)-\chi(K(B))$ of L3 for every finite $B$.

\begin{quote}
\textbf{Half-plane Jordan lemma.} Let $B$ be before $p$, with
$S\notin B$, $SE\in B$ and $W\in B$ or $SW\in B$. Then the 4-component of $S$ in
$\mathbb{Z}^2\setminus(B\cup\{p\})$ is finite if and only if $W$ or $SW$ lies in the
8-component of $SE$ in $B$.
\end{quote}

One direction says that an 8-path in $B$ from $L$ to $SE$, closed up through $p$,
cuts $S$ off from infinity. It has a Lean proof (above). The other says that if no
such path exists, $S$ escapes around the 8-component of $SE$
(\texttt{statement\_of\_enc\_join} derives \texttt{Statement} from this direction
alone). Both are instances of the digital Jordan curve theorem for (8,4)
connectivity~\cite{Rosenfeld1970,KongRosenfeld1989}. The second direction needs
a different tool, boundary following. Section~\ref{sec:crack} gives it. The
four-pixel diamond shows that the case cannot be avoided by a cleverer choice of
insertion order. Each of its pixels, inserted last in any of the eight scan
orders, falls into row 5 with $L$ and $SE$ joined.

The key statements, as they appear in the candidate file:
\begin{lstlisting}
theorem euler_insert_local (B : Finset Pixel) (p : Pixel) (hp : p ∉ B) :
    euler (insert p B) - euler B = euler (insert p (win B p)) - euler (win B p)

theorem C8_insert (hp : p ∉ B) :
    (C8 (insert p B) : ℤ) = C8 B + 1 - Nat.card {c // TouchC B p c}

theorem enclosed_insert (hp : p ∉ B)
    (hfin : Finite {c : (complementGraph B).ConnectedComponent // IsEnclosed c}) :
    (enclosed (insert p B) : ℤ) = enclosed B - (if IsEnclosed (cp B p hp) then 1 else 0)
      + Nat.card {x : {c : (complementGraph (insert p B)).ConnectedComponent // IsEnclosed c} //
          TouchPC B p x.1}

theorem statement_of_halfplane_jordan
    (hJ : ∀ (B : Finset Pixel) (p : Pixel), p ∉ B → Before B p →
      (p.1, p.2 - 1) ∉ B → (p.1 + 1, p.2 - 1) ∈ B →
      ((p.1 - 1, p.2) ∈ B ∨ (p.1 - 1, p.2 - 1) ∈ B) →
      (SEnc B p ↔ JoinLR B p)) :
    Statement

theorem senc_of_join (hb : Before B p) (hS : (p.1, p.2 - 1) ∉ B) (hJ : JoinLR B p) :
    SEnc B p

theorem invariant_nonneg (B : Finset Pixel) : 0 ≤ invariant B

theorem statement_of_le
    (h : ∀ B : Finset Pixel, (enclosed B : ℤ) ≤ C8 B - euler B) : Statement

lemma cross_horiz (v a b : Pixel) (hadj : G8.Adj a b)
    (ha1 : ¬ (a.1 = v.1 ∧ a.2 = v.2)) (ha2 : ¬ (a.1 = v.1 + 1 ∧ a.2 = v.2))
    (hb1 : ¬ (b.1 = v.1 ∧ b.2 = v.2)) (hb2 : ¬ (b.1 = v.1 + 1 ∧ b.2 = v.2)) :
    (cross v a b + cross (v.1 + 1, v.2) a b
      - (if a.1 = v.1 + 1 ∧ v.2 < a.2 then 1 else 0) - (if b.1 = v.1 + 1 ∧ v.2 < b.2 then 1 else 0))
      % 2 = 0 := by
  have := g8_coords hadj
  unfold cross
  simp only
  split_ifs <;> omega

lemma walk_lemma (v : (B : Set Pixel)) (y : ((insert p B : Finset Pixel) : Set Pixel))
    (h : (barrierGraph (insert p B)).Reachable (incl B p v) y) :
    (∃ y0 : (B : Set Pixel), y.1 = y0.1 ∧ (barrierGraph B).Reachable v y0) ∨ Touch B p v := by
  rw [SimpleGraph.reachable_iff_reflTransGen] at h
  induction h with
  | refl => exact Or.inl ⟨v, rfl, SimpleGraph.Reachable.refl _⟩
  | @tail y z _ hyz ih =>
    rcases ih with ⟨y0, hy, hr⟩ | ht
    · rw [barrier_adj_iff] at hyz
      by_cases hz : z.1 = p
      · right
        exact ⟨y0, hr, by rw [← hz, ← hy]; exact hyz.symm⟩
      · have hzB : z.1 ∈ (B : Set Pixel) := by
          have := z.2
          simp only [Finset.coe_insert, Set.mem_insert_iff] at this
          exact this.resolve_left hz
        left
        refine ⟨⟨z.1, hzB⟩, rfl, hr.trans (SimpleGraph.Adj.reachable ?_)⟩
        rw [barrier_adj_iff]; rw [← hy]; exact hyz
    · exact Or.inr ht
\end{lstlisting}


\section{A proof of L2, part 2: escape along the crack boundary}\label{sec:crack}

This section reports the 2026-11-09 session. It proves the remaining direction of the
half-plane Jordan lemma. With it, the file \texttt{proofs/P1.Insertion.lean} contains
\texttt{Candidate.main : ViridisStatements.P1.Insertion.Statement}. The project's
checker \texttt{viridis-lean-check} ran all five of its stages (denylist, compilation,
statement match, axioms, kernel replay) on this file and passed it, with axioms
\texttt{propext}, \texttt{Classical.choice} and \texttt{Quot.sound}. That is a local
preview; the verifier's own report is the certification page.

\paragraph{The claim.} Let $B$ be before $p$, with $S\notin B$, $SE\in B$, and suppose that
neither $W$ nor $SW$ lies in the 8-component $C$ of $SE$ in $B$. We show that $S$ is joined to
$E=p+(1,0)$ by a 4-path in $\mathbb{Z}^2\setminus(B\cup\{p\})$. Since $E$ has a free
ray, the component of $S$ is infinite (\texttt{join\_of\_senc}, by contraposition). Two
facts about $C$ make the problem local:
\begin{itemize}
\item no pixel of $B\setminus C$ lies within Chebyshev distance $1$ of $C$, since it
  would be 8-adjacent to $C$ and hence in $C$. So every pixel near $C$ but not in $C$ is free,
  apart from $p$ itself;
\item the only pixel of $C$ within Chebyshev distance $1$ of $p$ is $SE$. The other
  candidates are $W$ and $SW$ (not in $C$ by assumption), $S$ (free), and pixels at or
  after $p$ in scan order (not in $B$). So $p$ touches $C$ only at the corner shared by $SE$.
\end{itemize}

\paragraph{Cracks and their successor.} A \emph{crack} of a finite set $C$ is a pair
$(c,d)$ with $c\in C$, $d\in\{(\pm1,0),(0,\pm1)\}$ and $c+d\notin C$: a unit edge of
$K(C)$ with $C$ on one side and the outside pixel $c+d$ on the other. Let
$t=\mathrm{rot}(d)=(-d_2,d_1)$. The successor of $(c,d)$ follows the (8,4) rule:
\[
  \mathrm{next}(c,d)=\begin{cases}
  (c+d+t,\,-t) & \text{if } c+d+t\in C \quad\text{(turn outward, diagonal step in } C),\\
  (c+t,\,d) & \text{else if } c+t\in C \quad\text{(straight on)},\\
  (c,\,t) & \text{otherwise}\quad\text{(turn around } c).
  \end{cases}
\]
We give Lean proofs that $\mathrm{next}$ maps cracks to cracks (\texttt{next\_mem}) and
that it has an explicit left inverse $\mathrm{prev}$ on cracks, given by the mirror rule
(\texttt{prev\_next}). So $\mathrm{next}$ is a permutation of the finite set of cracks. By
a pigeonhole argument (\texttt{orbit}) every crack $x_0$ has a least period $n>0$ with
$\mathrm{next}^n(x_0)=x_0$.

\paragraph{The orbit from $S$.} Take $x_0=(SE,(-1,0))$, whose outside pixel is $S$,
and $x_E=(SE,(0,1))$, whose outside pixel is $E$. A direct computation gives
$\mathrm{next}(x_E)=x_0$: the diagonal pixel is $p\notin C$ and the side pixel is
$S\notin C$, so the boundary turns around $SE$. Since
$\mathrm{next}(\mathrm{next}^{n-1}x_0)=x_0=\mathrm{next}(x_E)$, injectivity gives
$\mathrm{next}^{n-1}x_0=x_E$. Minimality of $n$ gives $\mathrm{next}^{k}x_0\ne x_E$ for
$k<n-1$. Neither fact uses any planar topology.

\paragraph{One step is one short free path.} Let $x=(c,d)$ be a crack other than $x_E$.
We give a Lean proof (\texttt{crack\_step}) that the outside pixels of $x$ and of
$\mathrm{next}(x)$ are joined by a free 4-path of length at most $2$. In the first case
of the rule the outside pixel is unchanged. In the second it moves to the 4-adjacent
pixel $c+d+t$. In the third it moves from $c+d$ to $c+t$ through the diagonal pixel
$c+d+t$. Every pixel involved is within distance $1$ of $C$ and not in $C$, so it is
free unless it is $p$. The outside pixels are 4-adjacent to $C$, and $p$ is not. The diagonal pixel of
the third case equals $p$ only if $c=SE$ and $d+t=(-1,1)$, that is, only for $x=x_E$. Each
of these local facts is linear integer arithmetic over the four directions and is
closed by \texttt{omega}.

Chaining the steps for $k=0,\dots,n-2$ gives a free 4-path from $S$ to $E$
(\texttt{crack\_escape}). This proves the remaining direction. Then
\texttt{statement\_of\_enc\_join} gives \texttt{Statement}, and
\texttt{bridge\_of\_halfplane\_jordan} gives the bridge equation for every finite raster
(\texttt{bridge\_eq}).

\paragraph{Why this was the right tool.} The earlier sessions considered two
routes: boundary following, and the lemma that the border of a finite 4-connected set
is 8-connected. Formalizing boundary following usually means proving that the traced
sequence closes up and visits the right pixels. The permutation view avoids both. Closing
up is pigeonhole plus injectivity, and the only global information needed is that the
step into $x_0$ comes from $x_E$. That step is the one that passes the corner at $p$. We
also tried, on paper, an induction that fills free pixels of the component of $S$ one
at a time. It fails on annular components: filling any pixel of an annulus around a
one-pixel island merges the island with the outer barrier.

The new definitions and the final theorems:
\begin{lstlisting}
def next (C : Finset Pixel) (x : Pixel × Pixel) : Pixel × Pixel :=
  if (x.1.1 + x.2.1 - x.2.2, x.1.2 + x.2.2 + x.2.1) ∈ C then
    ((x.1.1 + x.2.1 - x.2.2, x.1.2 + x.2.2 + x.2.1), (x.2.2, -x.2.1))
  else if (x.1.1 - x.2.2, x.1.2 + x.2.1) ∈ C then ((x.1.1 - x.2.2, x.1.2 + x.2.1), x.2)
  else (x.1, (-x.2.2, x.2.1))

lemma prev_next (C : Finset Pixel) (x : Pixel × Pixel) (hx : x ∈ cracks C) :
    prev C (next C x) = x

lemma orbit {α : Type*} (f g : α → α) (s : Finset α) (hs : ∀ x ∈ s, f x ∈ s)
    (hg : ∀ x ∈ s, g (f x) = x) (x0 : α) (h0 : x0 ∈ s) :
    ∃ n, 0 < n ∧ f^[n] x0 = x0 ∧ ∀ k, 0 < k → k < n → f^[k] x0 ≠ x0

theorem crack_escape (hc0 : c0 ∈ C)
    (hfree : ∀ c ∈ C, ∀ q : Pixel, q ∉ C → Near q c → q ≠ p → q ∉ B)
    (hp : p = (c0.1 - 1, c0.2 + 1))
    (hpC : ∀ c ∈ C, Near p c → c = c0)
    (hS : (c0.1 - 1, c0.2) ∉ C) (hE : (c0.1, c0.2 + 1) ∉ C) :
    FR B p (c0.1 - 1, c0.2) (c0.1, c0.2 + 1)

theorem join_of_senc (hp : p ∉ B) (hb : Before B p) (hS : (p.1, p.2 - 1) ∉ B)
    (hSE : (p.1 + 1, p.2 - 1) ∈ B) (hs : SEnc B p) : JoinLR B p

theorem bridge_eq (B : Finset Pixel) : (enclosed B : ℤ) = C8 B - euler B

theorem main : ViridisStatements.P1.Insertion.Statement :=
  statement_of_enc_join fun B p hp hb hS hSE _ hs => join_of_senc B p hp hb hS hSE hs
\end{lstlisting}
Here \texttt{FR B p a b} says that $a$ and $b$ are free in $B\cup\{p\}$ and joined by a
4-path there, and \texttt{Near q c} says that $q$ and $c$ are at Chebyshev distance at most $1$.

\section{A proof of L3 and what it discharges}\label{sec:l3}

This section reports the 2026-11-10 session on the sealed statement \texttt{P1.Bridge}.

\subsection{The statement}

The module \texttt{ViridisStatements.P1.Bridge} restates the definitions of
Listing~\ref{lst:defs}. They are character for character the same as those of
\texttt{P1.Insertion}, which we checked by a textual diff. The statement is
\begin{lstlisting}
def Statement : Prop :=
  ∀ B : Finset Pixel,
    {c : (complementGraph B).ConnectedComponent | IsEnclosed c}.Finite ∧
      (enclosed B : ℤ) = C8 B - euler B
\end{lstlisting}
The first clause is needed because \texttt{enclosed} is a \texttt{Nat.card}, which
would read $0$ on an infinite type.

\begin{theorem}[L3]\label{thm:l3}
For every finite $B\subset\mathbb{Z}^2$, the set of finite 4-components of
$\mathbb{Z}^2\setminus B$ is finite, and its cardinality is $C_8(B)-\chi(K(B))$.
\end{theorem}

\begin{proof}[Proof sketch]
Finiteness (\texttt{finite\_enclosed}, 2026-11-08). Choose $M$ with $|b_1|,|b_2|<M$ for every
$b\in B$. A free pixel outside the box $[-M,M]^2$, or inside it but on a row that $B$
does not meet on that side, has a free horizontal ray, so its component is infinite
(\texttt{enclosed\_in\_box}). Choosing one representative pixel per enclosed component
therefore gives an injection into the finite box.
The equation (\texttt{bridge\_eq}, 2026-11-09). Order the pixels of $B$ in raster-scan order
and insert them one at a time, starting from $\varnothing$. The quantity
$\mathrm{enc}-C_8+\chi$ is $0$ on $\varnothing$, and by L2 every insertion preserves it.
\end{proof}

In the Lean file, L2 is restated over the \texttt{P1.Bridge} definitions as
\texttt{InsStatement} and proved as \texttt{Candidate.insertion}. The only new step for the
sealed target is to put the two clauses together:
\begin{lstlisting}
theorem main : ViridisStatements.P1.Bridge.Statement := fun B =>
  ⟨finite_enclosed B, bridge_eq B⟩
\end{lstlisting}
Here \texttt{finite\_enclosed} is stated as \texttt{Finite \{c // IsEnclosed c\}}. Lean accepts it
for the \texttt{Set.Finite} clause without conversion, because the two are definitionally
equal.

\subsection{Exactly one unbounded component}

The bridge counts \emph{bounded} compartments. \texttt{FaceCount} counts all of them,
and the ``$+1$'' in ``circuit rank $+\,1$'' is the unbounded one. We therefore also give a Lean
proof of the following (auxiliary theorems \texttt{unique\_unbounded} and
\texttt{card\_components}, which are not part of the sealed statement).

\begin{theorem}\label{thm:total}
For every finite $B$, $\mathbb{Z}^2\setminus B$ has exactly one infinite 4-component.
Hence it has finitely many 4-components in all, and their number is
$C_8(B)-\chi(K(B))+1$.
\end{theorem}

\begin{proof}[Proof sketch]
Take $M$ as above and $N=|M|$. Every row $y$ with $|y|\ge M$ and every column $x$ with
$|x|\ge M$ is free of $B$. Any free pixel with a coordinate of absolute value
at least $M$ therefore reaches $(N,N)$ along one free row and one free column
(\texttt{row\_reach}, \texttt{col\_reach}). The component of $(N,N)$ is infinite by a free ray.
An infinite component cannot lie inside the box, so it contains such a pixel and equals
the component of $(N,N)$. The count follows by splitting the components into enclosed and
non-enclosed ones and applying Theorem~\ref{thm:l3}.
\end{proof}

\begin{lstlisting}[caption={The total compartment count (excerpt from \texttt{proofs/P1.Bridge.lean}).}]
theorem unique_unbounded : ∃ c : (complementGraph B).ConnectedComponent,
    ¬ IsEnclosed c ∧ ∀ c', ¬ IsEnclosed c' → c' = c

theorem card_components : Finite (complementGraph B).ConnectedComponent ∧
    (Nat.card (complementGraph B).ConnectedComponent : ℤ) = C8 B - euler B + 1
\end{lstlisting}

\subsection{Which assumption of \emph{Trees Cannot Cut} this discharges}

The Lean core of \emph{Trees Cannot Cut}~\cite{Hart2026} (\texttt{DualCorridor.lean}) carries
the topological step as an explicit hypothesis:
\begin{lstlisting}
structure FaceCount (V : Type*) [Fintype V] [DecidableEq V] where
  faces : (G : SimpleGraph V) → [DecidableRel G.Adj] → ℕ
  duality : ∀ (G : SimpleGraph V) [DecidableRel G.Adj],
      (faces G : ℤ) = circuitRank G + 1
\end{lstlisting}
The manuscript justifies it by Alexander duality: the number of components of the fuel
matrix $\Omega\setminus C$ equals $\operatorname{rank}H_1(C\cup\partial\Omega)+1$, if we count the outer
compartment. The present result relates to this hypothesis as follows.

\paragraph{Discharged (on rasters).} The \emph{duality step} itself. For a corridor or
fuel break given as a finite raster, with the (8,4) reading, Theorems~\ref{thm:l3}
and~\ref{thm:total} prove with no hypotheses that the fuel matrix $\mathbb{Z}^2\setminus B$
has exactly $C_8(B)-\chi(K(B))+1$ compartments, of which $C_8(B)-\chi(K(B))$ are bounded.
Here $C_8-\chi$ is computed from the raster by counting components and local
$2\times 2$ patterns (L1). No appeal to Alexander duality, to a planar embedding or to
any topological library remains. The ``trees cannot cut'' message also has a raster
form: a raster with $\chi(K(B))=C_8(B)$ encloses nothing.

\paragraph{Not discharged.} \texttt{FaceCount} as stated is a hypothesis about an
abstract \texttt{SimpleGraph} and its \emph{graph} circuit rank $|E|-|V|+c(G)$. Our
theorem does not mention any such graph. Three links remain unformalised.
\begin{enumerate}
\item \emph{Raster-to-graph correspondence.} One needs a map from corridor graphs to
  rasters (or back) under which $\operatorname{circuitRank}(G)=C_8(B)-\chi(K(B))$. The obvious
  candidate, the 8-adjacency graph of the barrier pixels, fails: the solid $2\times2$ block
  has circuit rank $3$ and $C_8-\chi=0$ (negative control~1). Natural candidates are
  the 1-skeleton of $K(B)$ after collapsing filled squares, or a rasterisation of $G$ in which
  each edge becomes a one-pixel-wide digital arc and distinct edges meet only at vertex
  blocks. We have not formalised either.
\item \emph{Homology.} We do not prove that $C_8(B)-\chi(K(B))=\operatorname{rank}H_1(K(B))$,
  nor that $C_8(B)$ is the number of path components of $K(B)\subset\mathbb{R}^2$. The theorem is
  purely combinatorial, and it was designed so that it would not need these facts.
\item \emph{The landscape boundary.} The manuscript works in a bounded landscape $\Omega$
  and anchors the corridor to $\partial\Omega$ (the \texttt{anchor} construction of
  \texttt{DualCorridor.lean}). Our complement is taken in all of $\mathbb{Z}^2$, with no
  window. A windowed version, in which the frame of a rectangle is added to the barrier, is
  a special case of Theorem~\ref{thm:l3} applied to $B\cup\partial\Omega$. Matching it to the
  anchored graph again needs item~1.
\end{enumerate}
So the downstream results of \emph{Trees Cannot Cut} (zero-cut lemma, bridge and
non-bridge deletion, disjoint support) remain conditional on \texttt{FaceCount} in their
graph form. What changes is the status of the step that \texttt{FaceCount} stands for. On
rasters, it is now a theorem about $C_8-\chi$, and the remaining gap is the graph-theoretic
item~1, not a topological one.

\section{Computational evidence for L2 and L3}\label{sec:evidence}

The script \texttt{code/bridge\_check.py} is independent of both Lean models. It
was rewritten for this revision, because the first draft's script is not part of
the inputs available to this session. It pads each raster by one pixel, labels the 4-components of the complement and
the 8-components of $B$ by depth-first search, and computes $V-E+F$ directly. A complement component counts as enclosed when
it does not touch the padded border.
It found no failures (\texttt{results/bridge\_check.json}):
\begin{itemize}
\item L3 on all $2^{16}=65{,}536$ subsets of a $4\times4$ box;
\item L3 on 3{,}000 random rasters of side up to 12 with random density (seed 20261104);
\item L2 on $10{,}496$ insertions: each of the 256 configurations of the eight
  neighbours of $p$, alone and in 40 random contexts at Chebyshev distance
  between 2 and 6 from $p$.
\end{itemize}
This is evidence that the proposed statements are true and that their
conventions are the intended ones. It is not a proof.

For the 2026-11-08 reduction, the script \texttt{code/core\_check.py} (standard
library only, about 3 seconds) tests the exact hypothesis of
\texttt{statement\_of\_halfplane\_jordan}. It draws 4{,}000 random rasters (side 3
to 9, seed 20261108). At every candidate pixel $p$ it restricts the raster to the
pixels before $p$ and keeps the configurations of row 5 of
Table~\ref{tab:scan}. In all 17{,}344 such configurations (7{,}924 enclosed, 9{,}420
not enclosed), the finiteness of the component of $S$ agrees with the
8-connection of $L$ and $SE$. The script also checks the full invariant on 15{,}281
scan-order insertions, with no failures (\texttt{results/core\_check.json}). As a
mutation test, a variant that ignores $SW$ in the connection test reports 3{,}662
disagreements. So the check can tell the two sides apart.

\section{Limitations}

\begin{itemize}
\item The theorem concerns the combinatorial counts $V-E+F$ as defined in the
  sealed module. The development does not show that this number is the Euler
  characteristic of the topological space $K(B)\subset\mathbb{R}^2$. That
  identification is standard but is not formalised here, and L2/L3 are
  designed not to need it.
\item Only the 8-connected weighting of the bit quads is treated. The
  4-connected variant, with $+2Q_D$, is not formalised.
\item We give a Lean proof of L2 and of the sealed L3 statement \texttt{P1.Bridge}.
  The \texttt{FaceCount} hypothesis of \emph{Trees Cannot Cut} is about a corridor graph,
  not a raster. Discharging it also needs a raster-to-graph correspondence, the
  identification with homology and the boundary anchoring
  (Section~\ref{sec:l3}). None of these is formalised. The
  heuristic description of (J) in terms of arcs on the 8-cycle is a conjecture and
  has no Lean proof. The proof of L2 goes through scan order instead.
\item The bounded-box model used for the \texttt{decide} checks is not proved
  to agree with the semantic definitions. Only the four semantic control lemmas
  link the two.
\item The \texttt{decide} table of all 256 neighbourhood configurations
  uses only two global contexts: none, and a closed surrounding ring. The Python
  check adds 40 random contexts per configuration. Neither replaces a proof of
  L2.
\item Attempts to use the Aristotle prover (finiteness clause of L3 on
  2026-11-04; the full L2 statement on 2026-11-08; a status check on 2026-11-09)
  failed with an invalid API key, as in earlier sessions. On 2026-11-10 the prover was not
  needed and was not used. No automated proof search
  was done, and every proof in the file was written by hand.
\item L3 relates enclosed components to $C_8-\chi$, not to the circuit rank of
  a corridor graph. A paper that uses L3 in place of \texttt{FaceCount} must
  state which raster-to-graph correspondence it uses. The $2\times2$ block
  shows that the two notions differ.
\end{itemize}

\section{Reproducibility}
{\sloppy

The L1 proof file \texttt{proofs/P1.BitQuad.lean} and the script
\texttt{code/bitquad\_check.py} belong to the package of the 2026-10-28 run. It imports only
\texttt{Mathlib} and the sealed module \texttt{ViridisStatements.P1.BitQuad}
(sha256 \texttt{cd3cdeba48f6\ldots{}8e23}), and it builds with toolchain
\texttt{leanprover/lean4:v4.28.0} and Mathlib commit
\texttt{8f9d9cff6bd728b17a24e163c9402775d9e6a365}. The check takes about 20
seconds. The brute-force script runs with the Python~3 standard library:
\texttt{python3 code/bitquad\_check.py results/bitquad\_4x4.json}. No external
data are used.

The proposals \texttt{proposals/P1.Insertion.lean} and
\texttt{proposals/P1.Bridge.lean} each compile with \texttt{lean} against the
same toolchain and Mathlib commit, in about 100 and 60 seconds respectively,
including all \texttt{decide} checks (among them the two 256-configuration
tables) and the semantic control proofs. Both import only \texttt{Mathlib}. The evidence
for L2/L3 is reproduced by \texttt{python3 -I code/bridge\_check.py
results/bridge\_check.json} (standard library only, about 6 seconds). The
proposals, the evidence script and its results belong to the package of the
2026-11-04 run.

The file \texttt{proofs/P1.Insertion.lean} (2026-11-09 package) imports only
\texttt{Mathlib} and the sealed module \texttt{ViridisStatements.P1.Insertion}
(sha256 \texttt{800bd9b9fec5\ldots{}c734}). It builds with toolchain
\texttt{leanprover/lean4:v4.28.0} and the Mathlib commit above, and compiles in about
30 seconds with \texttt{lean}. Run
\texttt{viridis-lean-check proofs/P1.Insertion.lean P1.Insertion} to check it
(about 50 seconds). On our run it passed all five stages and reported the axioms
\texttt{propext}, \texttt{Classical.choice} and \texttt{Quot.sound}. The script
\texttt{code/core\_check.py} belongs to the 2026-11-08 session and is not part of
this package. The proof does not depend on it.

The file \texttt{proofs/P1.Bridge.lean} (2026-11-10 package, about 2{,}090 lines) imports only
\texttt{Mathlib} and the sealed module \texttt{ViridisStatements.P1.Bridge}
(sha256 \texttt{66038622a819\ldots{}171d}). It contains the whole development of L2
over the \texttt{P1.Bridge} definitions, so it does not depend on any earlier package.
It compiles in about 20 seconds with \texttt{lean} against the toolchain and Mathlib commit
above. Run \texttt{viridis-lean-check proofs/P1.Bridge.lean P1.Bridge} to check it
(about 40 seconds). On our run it passed all five stages, and \texttt{Candidate.main} reported the
axioms \texttt{propext}, \texttt{Classical.choice} and \texttt{Quot.sound}. The auxiliary
theorems \texttt{insertion}, \texttt{unique\_unbounded} and \texttt{card\_components}
depend on the same three axioms (\texttt{\#print axioms} in a scratch copy).

}

\section*{Changelog}
\begin{itemize}
\item 2026-10-28 (P1-T001): first version of the paper. Lean proof of L1
  (\texttt{P1.BitQuad}), proof architecture, brute-force check on the
  $4\times4$ box, and discussion of the role of L1 in L2/L3.
\item 2026-11-03 (P1-T002): new title and abstract. Added the proposed
  statements \texttt{P1.Insertion} (L2) and \texttt{P1.Bridge} (L3) with their
  conventions, controls (box-model \texttt{decide} checks and Lean proofs of the
  control values for the semantic definitions), brute-force evidence for L2/L3,
  the reference to Rosenfeld (1970), and updated limitations and reproducibility.
\item 2026-11-04 (P1-T002, second attempt): the verifier had rejected the first
  draft of the proposals because the import of the L1 module did not resolve in its
  environment. The proposal modules now restate the L1 definitions verbatim and
  import only Mathlib. Added the bitmask model with \texttt{decide} checks of all
  256 neighbourhood configurations in two contexts, and rewrote the evidence
  script (new seed). Updated abstract, limitations and reproducibility.
\item 2026-11-08 (P1-T003): new section on a reduction of L2
  (Section~\ref{sec:reduction}): Lean proofs of locality of $\Delta\chi$,
  of $\Delta C_8 = 1-k$, of finiteness of the enclosed components, of the change
  in the enclosed count, of L2 from the Jordan-type identity (J), and of L2 in
  the special case where the four 4-neighbours of $p$ are barrier. Then the
  raster-scan reduction: L2 from scan-order insertions, the explicit case analysis
  of Table~\ref{tab:scan}, and L2 from a half-plane digital Jordan lemma, which is
  checked by brute force (\texttt{code/core\_check.py}). A Lean proof of one
  direction of that lemma by crossing parity, and of the inequality
  $\mathrm{enc}\ge C_8-\chi$ for every finite raster. Updated the title, abstract, introduction,
  limitations and reproducibility. L2 remains unproved.
\item 2026-11-09 (P1-T003): new Section~\ref{sec:crack}. Lean proof of the remaining
  direction of the half-plane Jordan lemma by following the crack boundary
  (a permutation of the finite set of cracks), and hence a Lean proof of the sealed
  statement \texttt{P1.Insertion} (\texttt{Candidate.main}) and of the bridge
  equation $\mathrm{enc}=C_8-\chi$ (\texttt{bridge\_eq}). Updated the title, abstract,
  introduction, Table~\ref{tab:scan}, limitations and reproducibility.
\item 2026-11-10 (P1-T004): new Section~\ref{sec:l3}. Lean proof of the sealed L3
  statement \texttt{P1.Bridge} (\texttt{Candidate.main}), obtained from the 2026-11-09 development
  restated over the \texttt{P1.Bridge} definitions. New Lean proofs of uniqueness of the
  unbounded component and of the total compartment count $C_8-\chi+1$. A precise account
  of which part of \texttt{FaceCount} is discharged and which is not. Updated the title, abstract,
  introduction, limitations and reproducibility.
\end{itemize}

\bibliographystyle{plainnat}
\bibliography{refs}

\end{document}
